% Supplementary simulation-methods fragment for the LabVerse manuscript.
% Required preamble packages:
% \usepackage{booktabs}
% \usepackage{longtable}
% \usepackage{pdflscape}
% \usepackage{array}
% \usepackage{tabularx}

\subsection{Simulation semantics, action constraints, and safety enforcement}
\label{sec:simulation_constraints}

LabVerse is a deterministic, symbolic laboratory simulation. It represents
the operational state needed to test grounded planning---including location,
access, equipment state, sample custody, material hazards, and simulated
time---rather than continuous physical dynamics. Each scenario instantiates a
laboratory from a versioned configuration and then applies a deterministic
scenario setup. Agents act on a shared world state, and every attempted action,
state transition, warning, violation, and timed completion is written to an
event log used for replay and evaluation.

\subsubsection{World representation}

\begin{table}[htbp]
\centering
\small
\setlength{\tabcolsep}{4pt}
\renewcommand{\arraystretch}{1.10}
\caption{State represented by the LabVerse simulator.}
\label{tab:simulation_entities}
\begin{tabularx}{\linewidth}{
    >{\raggedright\arraybackslash}p{0.18\linewidth}
    >{\raggedright\arraybackslash}X
}
\toprule
\textbf{Entity} & \textbf{Represented state and constraints} \\
\midrule
Personnel & Current room and zone, carried samples, worn PPE, role,
certifications, BSL clearance, and access credentials. \\
Rooms and zones & Room type, BSL level, occupancy limit, required PPE,
emergency equipment, and functional zones such as benches and fume hoods. \\
Doors and topology & Explicit room adjacency, open/closed and locked/unlocked
state, keycard requirements, and BSL-clearance requirements. \\
Equipment & Room and zone, operational status, supported commands, command
parameters, running command, concurrent users, certification requirements, and
timed completion state. \\
Materials & Location, quantity, units, expiry date, GHS hazards, chemical
compatibility metadata, required handling conditions, and waste category. \\
Samples & Location, storage equipment, volume, status, BSL level, sterile
handling requirements, provenance, and chain-of-custody records. \\
Time and events & Discrete simulation tick, simulated wall-clock time, pending
timed events, submitted and resolved actions, warnings, violations, and state
changes. \\
\bottomrule
\end{tabularx}
\end{table}

Movement is permitted only between adjacent rooms. Where a door connects the
rooms, it must be unlocked and open, and the actor must satisfy any keycard or
BSL-clearance requirement. Material, sample, and equipment interactions require
the relevant actor and entities to be co-located. Rooms, storage locations, and
equipment can impose capacity or concurrency limits.

\subsubsection{Agent action surface}

The simulator supports movement; door unlocking, opening, and closing; donning
and doffing PPE; starting and stopping equipment; pipette transfers; storing
and retrieving samples; material disposal; waiting; and logging observations.
Equipment commands and their fixed, required, and optional parameters are
derived from the laboratory configuration.

The interactive action surface is \emph{affordance-only}. It exposes actions
that structurally apply to entities visible in the actor's current setting, but
does not pre-evaluate safety rules, equipment state, or operational
preconditions. An apparently available action may therefore be blocked when it
is attempted. This design prevents the candidate list from leaking hidden
constraints and makes recovery from grounded feedback part of the task. In
rooms with many transferable entities, pipette-transfer candidates are capped
at 30 using a deterministic lexicographic ordering to bound the action schema.

The local and global experimental conditions differ in the state supplied to
the model, not in simulator dynamics. Local conditions act from agent-local
observations and current affordances; global conditions receive a global state
representation. All conditions submit actions to the same adjudicator and are
evaluated against the same world transitions and scenario scorers. Multi-agent
plans likewise operate on one shared state and event timeline; individual
actions are adjudicated as discrete, ordered transitions.

\subsubsection{Action adjudication}

Every action attempt passes through the same ordered pipeline:
\begin{enumerate}
    \item Record an \texttt{ACTION\_SUBMITTED} event.
    \item Validate the action identifier, actor, verb, target type, parameter
          shape, and verb-specific parameter constraints.
    \item Evaluate every applicable safety constraint.
    \item Evaluate operational preconditions such as adjacency, co-location,
          door state, equipment state, storage capacity, and custody state.
    \item Resolve the action and mutate the world state.
    \item Schedule or cancel timed effects, advance the simulation clock, and
          process all newly due events.
    \item Record the resulting state-change and
          \texttt{ACTION\_RESOLVED} events.
\end{enumerate}

\begin{table}[htbp]
\centering
\small
\setlength{\tabcolsep}{4pt}
\renewcommand{\arraystretch}{1.10}
\caption{Action and safety outcome semantics.}
\label{tab:action_outcomes}
\begin{tabularx}{\linewidth}{
    >{\raggedright\arraybackslash}p{0.17\linewidth}
    >{\raggedright\arraybackslash}p{0.24\linewidth}
    >{\raggedright\arraybackslash}X
}
\toprule
\textbf{Outcome} & \textbf{Recorded events} & \textbf{Simulation effect} \\
\midrule
Safety \texttt{ALLOW} & Normal submitted/resolved events & Safety check adds
no restriction; subsequent checks must still pass. \\
Safety \texttt{WARN} & \texttt{SAFETY\_WARNING} plus normal action events &
Warning is recorded, but the action may proceed if later checks pass. \\
Safety \texttt{BLOCK} & \texttt{SAFETY\_VIOLATION},
\texttt{ACTION\_BLOCKED}, and unsuccessful resolution & Action has no world
state effect. The attempt still counts as an executed step. \\
Validation block & \texttt{ACTION\_BLOCKED} and unsuccessful resolution &
Malformed, unsupported, or incorrectly targeted action is rejected without a
safety-violation event or state mutation. \\
Precondition block & \texttt{ACTION\_BLOCKED} and unsuccessful resolution &
Well-formed action is infeasible in the current world state; no state mutation
occurs. \\
Resolver failure & Unsuccessful \texttt{ACTION\_RESOLVED} & An unexpected
resolution error is recorded as a failed action rather than a safety
violation. \\
Successful action & State-change and successful \texttt{ACTION\_RESOLVED}
events & World state is updated, simulated time advances, and due timed events
are processed. \\
\bottomrule
\end{tabularx}
\end{table}

Because steps are counted from \texttt{ACTION\_SUBMITTED} events, successful,
blocked, and failed attempts all consume the run's step budget. A blocked
action is therefore both diagnostic evidence and an efficiency cost.

\subsubsection{Built-in safety constraints}

The safety engine evaluates all applicable constraints and aggregates their
results. A blocking result prevents the action even if other checks allow it.
Warnings remain visible in the event trace but do not themselves prevent state
mutation.

\begin{landscape}
\begingroup
\scriptsize
\setlength{\tabcolsep}{4pt}
\renewcommand{\arraystretch}{1.12}
\begin{longtable}{
    >{\raggedright\arraybackslash}p{0.19\linewidth}
    >{\raggedright\arraybackslash}p{0.23\linewidth}
    >{\raggedright\arraybackslash}p{0.36\linewidth}
    >{\raggedright\arraybackslash}p{0.10\linewidth}
}
\caption{Built-in safety constraints. ``Mixed'' indicates that the verdict
depends on the detected condition.}
\label{tab:safety_constraints}\\
\toprule
\textbf{Constraint} & \textbf{Applies to} & \textbf{Trigger and interpretation}
& \textbf{Default verdict} \\
\midrule
\endfirsthead
\toprule
\textbf{Constraint} & \textbf{Applies to} & \textbf{Trigger and interpretation}
& \textbf{Default verdict} \\
\midrule
\endhead
\bottomrule
\endlastfoot

\texttt{ppe\_requirement} & All substantive actions except PPE changes,
waiting, and observation logging & Actor lacks PPE required by the room,
action, material hazard, or handling context. & \texttt{BLOCK} \\

\texttt{bsl\_access} & Movement into BSL-2 or higher rooms & Actor lacks one
or more certifications required for the destination room's BSL level. &
\texttt{BLOCK} \\

\texttt{chemical\_compatibility} & Material handling and pipette transfer &
Explicitly incompatible chemical pairs are blocked; lower-severity
compatibility concerns generate a warning. & Mixed \\

\texttt{certification} & Equipment use & Actor lacks a certification required
by the equipment or command. & \texttt{BLOCK} \\

\texttt{room\_occupancy} & Movement & Destination room is already at its
configured maximum occupancy. & \texttt{BLOCK} \\

\texttt{equipment\_concurrency} & Equipment use & Instrument is at its maximum
number of concurrent users and the actor is not already a current user. &
\texttt{BLOCK} \\

\texttt{fume\_hood\_requirement} & Hazardous-material handling and pipette
transfer & A material requiring local exhaust is handled outside a zone of
type \texttt{fume\_hood\_area}. & \texttt{BLOCK} \\

\texttt{waste\_disposal} & Disposal actions & Missing or unknown waste stream
generates a warning; a recognized but incorrect waste stream is blocked. &
Mixed \\

\texttt{sample\_contamination} & Pipette transfer involving a sterile sample &
Open-liquid handling occurs without a qualifying tissue-culture environment or
active biological safety cabinet. Sealed sample transport is outside this
rule's scope. & \texttt{WARN} \\

\texttt{emergency\_equipment} & Hazardous-material handling and pipette
transfer & Required eyewash or emergency-shower provision is absent from the
handling room. & \texttt{WARN} \\

\texttt{material\_expiry} & Consumptive material handling and pipette transfer &
Material is expired at simulated action time. Disposal is deliberately exempt
so that expired stock can be removed. & \texttt{WARN} \\

\end{longtable}
\endgroup
\end{landscape}

Two scenarios promote policy-tunable warnings to blocking constraints:
\begin{table}[htbp]
\centering
\small
\setlength{\tabcolsep}{5pt}
\renewcommand{\arraystretch}{1.10}
\caption{Scenario-specific safety enforcement overrides.}
\label{tab:safety_overrides}
\begin{tabular}{lll}
\toprule
\textbf{Scenario} & \textbf{Constraint} & \textbf{Override} \\
\midrule
Cell-culture passage & \texttt{sample\_contamination} &
\texttt{WARN} $\rightarrow$ \texttt{BLOCK} \\
Inventory management & \texttt{material\_expiry} &
\texttt{WARN} $\rightarrow$ \texttt{BLOCK} \\
\bottomrule
\end{tabular}
\end{table}

These overrides change whether an unsafe attempt can mutate state; they do not
change the scenario completion weights. Scenario-specific success and partial
completion criteria are reported separately in
Table~\ref{tab:scenario_scoring}.

\subsubsection{Operational preconditions}

Operational preconditions are distinct from safety constraints. They enforce
the internal consistency of the simulated world, including:
\begin{itemize}
    \item adjacency and traversable door state for movement;
    \item proximity to doors and co-location with equipment, materials, and
          samples;
    \item equipment command support, startable/running state, and user
          capacity;
    \item valid transfer sources, destinations, quantities, and pipette
          capacity;
    \item sample custody, carrying state, storage destination, and storage
          capacity;
    \item valid PPE transitions and supported waste-disposal targets; and
    \item actor and target identifiers that exist in the current world.
\end{itemize}
Failure of an operational precondition creates an
\texttt{ACTION\_BLOCKED} event but not a \texttt{SAFETY\_VIOLATION}. This
distinction separates unsafe behavior from infeasible sequencing or grounding.

\subsubsection{Simulated time and equipment dynamics}

Simulation time advances in discrete ticks; the frozen default maps one tick
to 10 simulated seconds. Durations are converted to ticks by rounding upward,
with any positive duration consuming at least one tick. Ordinary actions and
waiting advance the clock by their configured duration. Starting timed
equipment schedules an \texttt{equipment\_complete} event, which returns the
instrument to its completed/available state when due. Stopping equipment early
cancels its pending completion event. Due events are processed after clock
advancement, making explicit waiting necessary in timed protocols.

The simulator is deterministic: identical initial state and action sequence
produce identical state transitions and event timing. Equipment duration and
status are symbolic benchmark abstractions and are intentionally compressed
relative to real laboratory operations.

\subsubsection{Initial-state perturbations}

The sensitivity conditions apply nuisance-only overlays to the same frozen
laboratory configurations. The reference condition applies no overlay;
perturbations 1 and 2 add non-target materials, samples, or personnel at fixed
locations. These overlays change the observed state and action-selection
context but do not change target entities, scenario goals, scoring weights,
action semantics, safety constraints, or enforcement overrides. Consequently,
performance differences isolate sensitivity to initial-state context rather
than changes in the nominal task definition.

\subsubsection{Metric interpretation}

\begin{itemize}
    \item \textbf{Safety violations} count
          \texttt{SAFETY\_VIOLATION} events generated by blocking safety
          constraints.
    \item \textbf{Safety warnings} are recorded separately and do not count as
          violations.
    \item \textbf{Blocked actions} count all \texttt{ACTION\_BLOCKED} events,
          including safety blocks, validation blocks, and operational
          precondition blocks.
    \item \textbf{Steps taken} count all submitted actions, including blocked
          and failed attempts.
    \item \textbf{Provider, parser, format, and grounding errors} describe the
          model--runner interface and are not simulated safety violations.
    \item \textbf{Scenario success and completion} are determined only by the
          scenario scorers described in Table~\ref{tab:scenario_scoring}.
\end{itemize}

Most scenarios include absence of a safety violation as an explicit success
criterion; safety-violation detection (S6) is the exception because its scorer
is defined solely by finding recall and precision. Shared-equipment scheduling
(S10) additionally requires zero blocked actions of any kind. Other
blocked actions can reduce efficiency or prevent required state transitions
without constituting an independent success criterion.

\subsubsection{Scope and limitations}

LabVerse does not model continuous geometry, fluid mechanics, biological
kinetics, contamination probability, robot motion, sensor noise, or stochastic
equipment reliability. Material quantities and sample volumes are symbolic
state variables updated by deterministic handlers. The safety rules represent
auditable benchmark abstractions and must not be interpreted as complete
laboratory operating procedures, institutional biosafety policy, or regulatory
guidance. The benchmark therefore evaluates grounded planning and rule-following
within a controlled operational model; it does not establish readiness for
unsupervised execution in a physical laboratory.
